% Explicit consequences of the declared 2013 HIT principles, not MLTT existence proofs.
\section{Higher inductive types}

\begin{annexitem}{annex:p07-hit-path-calculus}{thm:p07-hit-path-calculus}{Typed transport and dependent composition}
\begin{proof}
Define \(c(p,u)\) by identity elimination on \(p\), taking reflexivity
at reflexivity. The constant-family comparison for \(\operatorname{apd}\)
is also proved by eliminating \(p\); both sides then become reflexivity.
Eliminating \(p\) in the equality-family formula reduces it to
\(r=\mathrm{refl}^{-1}\cdot r\cdot\mathrm{refl}\), proved by the unit
paths. The fixed-endpoint formula is the same calculation without the last
factor. Transport composition \(\gamma_P\) is defined by elimination on
\(p,q\), with reflexivity at two reflexivities. Concatenating it with the
action on \(h\) and then \(k\) gives exactly a path from
\(\operatorname{tr}_P(p\cdot q,u)\) to \(w\). Elimination on
\(p,q\) proves \(\mu_f\), again at reflexivity. For \(\alpha:p=q\),
\(\operatorname{tr}^2\) is ordinary action on that equality;
elimination on \(\alpha\) defines \(\operatorname{apd}^2\), whose
reflexivity case is the unit comparison of \(\operatorname{apd}_f(p)\)
with \(\mathrm{refl}\cdot\operatorname{apd}_f(p)\).
Every unit, inverse, associativity and action-on-concatenation/inversion
comparison used below is obtained by identity elimination of its paths,
with reflexivity in the reflexivity case. Cancellation composes on the left
with the inverse common factor, then uses those associativity/inverse/unit
paths. No such comparison is asserted judgmental at an arbitrary path.
\end{proof}
\end{annexitem}

\begin{annexitem}{annex:p07-hit-circle-recursion}{thm:p07-hit-circle-recursion}{The constant-family computation conversion}
\begin{proof}
In \(P(x):=B\), supply \(c(\mathrm{loop},b)\cdot\ell\) to the
dependent eliminator. It gives \(r(\mathrm{base})\equiv_{\mathrm j}b\)
and \(\operatorname{apd}_r(\mathrm{loop})=c(\mathrm{loop},b)\cdot\ell\).
The preceding lemma identifies the left side with
\(c(\mathrm{loop},r(\mathrm{base}))\cdot\operatorname{ap}_r(\mathrm{loop})\).
The two \(c\)-factors are identical by judgmental point computation.
Cancel that factor by the proved path operations. This gives the asserted
propositional computation \(\operatorname{ap}_r(\mathrm{loop})=\ell\).
\end{proof}
\end{annexitem}

\begin{annexitem}{annex:p07-hit-circle-universal}{thm:p07-hit-circle-universal}{Coherent uniqueness and both inverse homotopies}
\begin{proof}
Induct in \(P(x):=f(x)=g(x)\), starting at the stated \(p\).
Its transport around the loop is
\(\operatorname{ap}_f(\mathrm{loop})^{-1}\cdot p\cdot
\operatorname{ap}_g(\mathrm{loop})\). The specified coherence, composed
on the left with the inverse first factor, identifies this with \(p\).
It is the needed dependent lift, so induction gives a homotopy and funext
gives a function path.

Let \(E\) be evaluation and \(R\) recursion. In \(E(R(b,\ell))\),
the point is judgmentally \(b\) and the loop component has the specified
path to \(\ell\). A pair path gives \(E(R(b,\ell))=(b,\ell)\):
pair paths over \((p,q)\), with \(q:\operatorname{tr}(p,z)=z'\),
are constructed by eliminating \(p\) and applying the pair constructor
to \(q\). Here the first component is reflexivity.
For \(f\), the map \(r=R(E(f))\) agrees at the basepoint judgmentally,
and its path computation identifies \(\operatorname{ap}_r(\mathrm{loop})\)
with \(\operatorname{ap}_f(\mathrm{loop})\). The uniqueness argument with
point path reflexivity gives \(r\sim f\); funext gives \(R(E(f))=f\).
Thus both inverse homotopies are supplied. Their conversion to the asserted
fiber-equivalence uses the complete adjointification/fiber proofs
\cref{thm:p07-hott-adjointification,thm:p07-hott-fiber-equivalence}; it does
not infer fiber propositionality from raw uniqueness alone.
\end{proof}
\end{annexitem}

\begin{annexitem}{annex:p07-hit-interval-contractible}{thm:p07-hit-interval-contractible}{The dependent segment coherence}
\begin{proof}
In \(P(x):=x=1_I\), choose \(\mathrm{seg}\) at \(0_I\) and reflexivity
at \(1_I\). The segment lift asks that transport of \(\mathrm{seg}\)
along itself be reflexivity. The fixed-endpoint transport comparison gives
\(\mathrm{seg}^{-1}\cdot\mathrm{seg}\), and the inverse law supplies
the lift. Induction gives every \(x=1_I\); invert these paths for a
contraction with center \(1_I\) in the usual center-to-point orientation.
\end{proof}
\end{annexitem}

\begin{annexitem}{annex:p07-hit-interval-extensionality}{thm:p07-hit-interval-extensionality}{A function path without invoking funext}
\begin{proof}
For each \(a\), interval recursion gives \(k_a:I\to B\) with
endpoints \(f(a),g(a)\) and segment image \(h(a)\).
Put \(K(i):=\lambda a.k_a(i)\). Point computation and the declared
judgmental function \(\eta\) give \(K(0_I)\equiv_{\mathrm j}f\) and
\(K(1_I)\equiv_{\mathrm j}g\). Hence \(\operatorname{ap}_K(\mathrm{seg})\)
is a path \(f=g\). No funext step is used, and no additional computation
rule for applying this path to each \(a\) is asserted.
\end{proof}
\end{annexitem}

\begin{annexitem}{annex:p07-hit-circle-nontrivial}{thm:p07-hit-circle-nontrivial}{Univalent boolean monodromy}
\begin{proof}
The boolean swap \(e:\mathbf2\simeq\mathbf2\) is its own inverse;
both inverse homotopies are the two boolean cases. Circle recursion into
the declared universe gives \(C\) with base fiber \(\mathbf2\) and
\(\operatorname{ap}_C(\mathrm{loop})=\operatorname{ua}(e)\).
Identity elimination proves
\(\operatorname{tr}_C(p,u)=
\operatorname{tr}_{X\mapsto X}(\operatorname{ap}_C(p),u)\), since both
sides are \(u\) at reflexivity. The propositional univalence computation
of \cref{def:p07-hott-univalence} consequently gives
\(\operatorname{tr}_C(\mathrm{loop},\mathrm{false})=\mathrm{true}\).
If the loop were reflexivity, action on that equality would make this
transport \(\mathrm{false}\), contradicting boolean no confusion.
For that last fact, define a boolean family with value \(\mathbf1\) at
false and \(\mathbf0\) at true; transporting its unit inhabitant along
a proposed false-to-true path would inhabit \(\mathbf0\).
If the circle were a set, its two parallel basepoint loops would be equal,
again a contradiction. Nothing here identifies all circle loops with integers.
\end{proof}
\end{annexitem}

\begin{annexitem}{annex:p07-hit-suspension-circle}{thm:p07-hit-suspension-circle}{Both meridian coherences}
\begin{proof}
Write \(m_0,m_1\) for the false and true meridians, and use the stated
recursors \(F,G\). For \(GF(x)=x\), suspension induction starts with
reflexivity at north and \(m_1\) at south. The transported north component
along \(m_i\) is \(\operatorname{ap}_{GF}(m_i)^{-1}\cdot m_i\).
For false it is
\((m_0\cdot m_1^{-1})^{-1}\cdot m_0
=m_1\cdot m_0^{-1}\cdot m_0=m_1\).
For true it is \(\mathrm{refl}^{-1}\cdot m_1=m_1\).
Boolean elimination gives both lifts, hence the full homotopy.
For \(FG(x)=x\), circle induction starts with reflexivity, and
\(\operatorname{ap}_{FG}(\mathrm{loop})
=\operatorname{ap}_F(m_0)\cdot\operatorname{ap}_F(m_1)^{-1}
=\mathrm{loop}\cdot\mathrm{refl}^{-1}=\mathrm{loop}\).
This gives the required equality-family loop lift. The resulting two
inverse homotopies yield equivalence by the cited adjointification/fiber
criterion, with all displayed computations propositional where appropriate.
\end{proof}
\end{annexitem}

\begin{annexitem}{annex:p07-hit-pushout-universal}{thm:p07-hit-pushout-universal}{The triple of cocone data}
\begin{proof}
Evaluation sends \(t\) to \((t\circ\mathrm{inl},t\circ\mathrm{inr},
\lambda c.\operatorname{ap}_t(\mathrm{glue}(c)))\). Recursion sends
\((i,j,h)\) to its prescribed point and glue values. Evaluation after
recursion has its first two components judgmentally \(i,j\), using
function \(\eta\). The third component equals \(h\) by funext applied
to the propositional glue computations. Pair paths over reflexivity give
the data inverse homotopy; the third component is not promoted to judgmental
equality. For \(t\), let \(r\) be recursion on its evaluated data.
In \(P(x):=r(x)=t(x)\), use reflexivity on both point constructors.
The transported glue component is
\(\operatorname{ap}_r(\mathrm{glue}(c))^{-1}\cdot
\operatorname{ap}_t(\mathrm{glue}(c))\); the computation witness and
inverse cancellation identify it with reflexivity. Dependent induction gives
\(r\sim t\), and funext the second inverse path. The established
adjointification/fiber theorem gives equivalence.
When both outside types are \(\mathbf1\), the constructors and full
dependent eliminator are the suspension ones. Maps between these two
presentations on the poles and meridians have inverse homotopies by this same
equality-family induction. With middle type \(\mathbf2\), suspension-circle
equivalence and the non-set result prove the failure to preserve sets.
\end{proof}
\end{annexitem}

\begin{annexitem}{annex:p07-hit-truncation-elimination}{thm:p07-hit-truncation-elimination}{All truncation lifts and inverse homotopies}
\begin{proof}
In a proposition-valued family, transported and target values at a squash
path belong to the same proposition, so its proposition witness gives
the required lift. In a set-valued family, a dependent 2-lift has type
\(h=\operatorname{tr}^2_P(\alpha,u)\cdot k\), where both sides are
parallel paths in \(P(y)\). Its set witness gives exactly this equality.
Thus the stated raw dependent principles supply the restricted eliminators.

We also need the familiar fact that a proposition is a set. For an inhabited
proposition with point \(a\), let \(H(x):a=x\) be its proposition
witness. Naturality, proved by identity elimination, gives
\(H(x)\cdot p=H(y)\) for \(p:x=y\). Consequently every such \(p\)
equals \(H(x)^{-1}\cdot H(y)\), so any two parallel paths are equal.
This proves the set assertion without a UIP axiom. Under funext, a function
type with proposition-valued codomain is also a proposition, by applying
the proposition witnesses pointwise.

For each universal property, the inverse to precomposition is the recursor.
Point computation and funext give precomposition after recursion equal to
the original map. In the propositional case, any two maps into \(B\)
are equal by the preceding pointwise argument. In the set case, for
\(h:\|A\|_0\to B\), the equality family comparing \(h\) with the
recursor of its point values is proposition-valued, hence set-valued.
Set-truncation induction extends its reflexivity point data to a homotopy;
funext gives the other inverse path. Adjointification/fiber equivalence
then gives both displayed equivalences.
\end{proof}
\end{annexitem}

\begin{annexitem}{annex:p07-hit-truncation-boundaries}{thm:p07-hit-truncation-boundaries}{Empty elimination and no arbitrary witness extraction}
\begin{proof}
The empty type is a proposition: its first argument eliminates any proposed
pair of points. Truncation recursion on \(\mathbf0\to\mathbf0\)
therefore gives \(\|\mathbf0\|\to\mathbf0\). Together with the empty
eliminator in the other direction, the inverse homotopies are vacuous, so
these maps give the asserted equivalence. A proposition with a point \(b\)
is contractible by its witness \(\prod_x(b=x)\). Use \(|\star|\) for
\(\|\mathbf1\|\) and \(|\mathrm{false}|\) for \(\|\mathbf2\|\).
If the stated \(e\) existed, action on
\(\mathrm{squash}(|\mathrm{false}|,|\mathrm{true}|)\), followed by its
two claimed point equalities, would give false equal to true. The boolean
family proof of no confusion above makes this impossible.
\end{proof}
\end{annexitem}

\begin{annexitem}{annex:p07-hit-quotient-universal}{thm:p07-hit-quotient-universal}{Mere representatives and coherent quotient recursion}
\begin{proof}
For mere surjectivity, induct in the proposition
\(P(x):=\|\sum_a(q(a)=x)\|\). At \(q(a)\), use
\(|(a,\mathrm{refl})|\). Proposition-valued fibers supply the glue and
truncation lifts just proved, so quotient induction gives all \(P(x)\).
This proof never eliminates that bracket into a nonproposition to choose
a representative.

Evaluate a map as its point values and \(\operatorname{ap}\) on each
relation constructor. Recursion from \((g,\rho)\) gives its inverse.
The first evaluated component computes to \(g\); the relation computations
give the second equality by funext. Pair paths give the data inverse
homotopy. For a quotient map \(h\), compare it with recursion on its
evaluated data. Their equality at any \(x\) is a proposition because
\(B\) is a set. Quotient dependent induction with reflexivity point data
therefore gives the other homotopy, with all lifts furnished by the
proposition-valued family. Funext and adjointification/fiber equivalence
complete the universal property. The final warning follows because ordinary
identity is always reflexive and symmetric, while a mere relation need not be.
\end{proof}
\end{annexitem}

\begin{annexitem}{annex:p07-hit-quotient-example}{thm:p07-hit-quotient-example}{The directed boolean relation}
\begin{proof}
Take center \(q(\mathrm{false})\). Boolean elimination supplies paths
from this center to each generator: reflexivity for false, and
\(\mathrm{eq}(\star)\) for true. Since \(Q\) is a set, the family of
paths from the center is proposition-valued. Quotient induction extends
those point paths to every \(x:Q\), with automatic higher lifts, giving
a contraction. The reflexivity path at the false image is inhabited,
whereas its displayed \(R\)-fiber is the empty type by definition.
\end{proof}
\end{annexitem}

\begin{annexitem}{annex:p07-hit-hub-contraction}{thm:p07-hit-hub-contraction}{Transport of the based-map witness}
\begin{proof}
Choose center \((c_b,\mathrm{refl}_b)\), where \(c_b\) is constant.
For \((r,p)\), put
\(L(x):=\mathrm{spoke}_r(x)\cdot\mathrm{spoke}_r(s_0)^{-1}\cdot p\).
Funext gives \(d:r=c_b\), with application at \(s_0\) identified
propositionally with \(L(s_0)\). The path laws identify \(L(s_0)\)
with \(p\). Transport in \(r\mapsto r(s_0)=b\) sends \(p\) to
\(\operatorname{ap}_{r\mapsto r(s_0)}(d)^{-1}\cdot p\); funext
computation and the preceding identification reduce this to
\(p^{-1}\cdot p=\mathrm{refl}_b\). The pair-path construction thus
gives \((r,p)=(c_b,\mathrm{refl}_b)\). Invert this pair path to obtain
the center-to-point contraction. No sphere--loop equivalence is assumed
in this particular contraction argument.
\end{proof}
\end{annexitem}

\begin{annexitem}{annex:p07-hit-hub-dependent}{thm:p07-hit-hub-dependent}{The actual supplied fiber contraction data}
\begin{proof}
For \(r:S^{n+1}\to\|A\|_n\) and
\(r':\prod_xP(r(x))\), define
\(t(x):=\operatorname{tr}_P(\mathrm{spoke}_r(x),r'(x))\) in the fixed
fiber \(P(\mathrm{hub}(r))\). Apply the assumed contraction function
for that fiber to \(t\), obtaining \(u\) and
\(v:\prod_x(t(x)=u)\). Set \(H(r,r'):=u\) and
\(S(r,r',x):=v(x)\). These have exactly the raw hub and dependent-spoke
types. Apply the declared eliminator with \(g,H,S\); its point
computation gives the conclusion. This proof constructs the missing raw
data from the explicit hypothesis and does not manufacture the general
\(n\)-type sphere-contraction theorem.
\end{proof}
\end{annexitem}
